\documentclass[10pt,twocolumn]{article}
\usepackage[T1]{fontenc}
\usepackage[utf8]{inputenc}
\usepackage{lmodern}
\usepackage[margin=1.8cm,top=2.4cm,bottom=2cm,columnsep=0.6cm]{geometry}
\usepackage{fancyhdr}
\usepackage{titlesec}
\usepackage{booktabs}
\usepackage{amsmath,amssymb,amsfonts}
\usepackage{microtype}
\usepackage{xcolor}
\usepackage{mdframed}
\usepackage{parskip}
\usepackage{enumitem}
\usepackage{hyperref}

\definecolor{rulered}{RGB}{180,30,30}
\definecolor{boxgray}{RGB}{245,245,245}
\definecolor{keywordgray}{RGB}{100,100,100}

\pagestyle{fancy}
\fancyhf{}
\fancyhead[L]{\textsc{\small Claude Mag}}
\fancyhead[C]{\small\textit{Machine Learning Research Digest}}
\fancyhead[R]{\small 2026-11-25}
\fancyfoot[C]{\small\thepage}
\renewcommand{\headrulewidth}{0.8pt}

\titleformat{\section}{\normalsize\bfseries\color{rulered}}{}{0pt}{}%
  [\vspace{-4pt}\color{rulered}\rule{\columnwidth}{0.4pt}]
\titlespacing{\section}{0pt}{8pt}{4pt}

\hypersetup{colorlinks=true,urlcolor=rulered,linkcolor=black}

\begin{document}
\setlength{\parindent}{0pt}
\setlength{\parskip}{4pt}

{\small\color{keywordgray}\textbf{Keywords:}
  SFT Corruption \textbullet{} RLVR \textbullet{} Reasoning Recovery \textbullet{} Form vs.\ Function \textbullet{} Data Quality}\par\vspace{4pt}

{\LARGE\bfseries\raggedright Reinforcement Learning Unlearns Corrupted SFT Reasoning}\par\vspace{4pt}
{\large\itshape\raggedright RL recovers near-full reasoning from corrupted SFT traces, revealing that SFT installs output format while RL independently discovers functional reasoning}\par\vspace{6pt}
{\small\color{keywordgray}arXiv:2611.26756 \textbullet{} cs.LG, cs.CL \textbullet{} submitted 19 November 2026}
\par\vspace{2pt}\noindent{\color{rulered}\rule{\columnwidth}{0.6pt}}\vspace{6pt}

A model SFT-trained on heavily corrupted chain-of-thought traces recovers near-full reasoning accuracy through RL fine-tuning.
SFT accuracy is a poor predictor of post-RL accuracy, suggesting that SFT installs structural output format while RL independently rediscovers functional reasoning---two separable contributions from two separable training stages.

\begin{mdframed}[backgroundcolor=boxgray,linecolor=rulered,linewidth=1pt,
    innerleftmargin=6pt,innerrightmargin=6pt,
    innertopmargin=6pt,innerbottommargin=6pt]
\textbf{\small AT A GLANCE}\par\vspace{4pt}
\begin{description}[leftmargin=0pt,labelsep=4pt,itemsep=2pt,parsep=0pt,
    font=\small\bfseries]
  \item[Problem:] The standard view holds that SFT trace quality determines post-RL performance; this motivates expensive chain-of-thought curation.
  \item[Why it matters:] If SFT trace content is not load-bearing, teams are over-investing in data curation and misidentifying training failure modes.
  \item[Method:] Five corruption levels on GSM8K SFT data; measure pre-RL and post-RL accuracy; diagnostic attention ablations.
  \item[Result:] Post-RL accuracy stays within 12 points of clean-data performance even at full SFT corruption; SFT$\to$post-RL correlation is $r = 0.41$.
  \item[Evidence:] GSM8K; GRPO; binary correctness reward; five corruption levels from 20\% step replacement to full noise.
\end{description}
\end{mdframed}

\section{The Big Picture}

The two-stage SFT-then-RL recipe for training reasoning LLMs is widely used, but the relative contributions of each stage are poorly understood.
The standard assumption is that SFT quality determines the ceiling on post-RL performance: a model SFT-trained on correct, high-quality chain-of-thought traces starts RL from a better initialization than one trained on noisy or incorrect traces.

This assumption has shaped significant research investment.
Teams expend substantial effort curating, filtering, and distilling high-quality SFT data for reasoning tasks, on the grounds that every quality improvement in the SFT set propagates to post-RL performance.

This paper tests the assumption with a controlled corruption study and finds it largely wrong.
RL is a powerful enough optimizer to recover most reasoning ability from a poor initialization, as long as the SFT stage has established the structural scaffolding of step-by-step responses.

\section{Why Existing Approaches Fall Short}

\textbf{SFT-focused data curation} has been the dominant response to reasoning failure.
When a trained model performs poorly, the first diagnostic is typically the SFT dataset---is it large enough, is it clean, does it cover the right distribution?
This paper suggests the RL stage should be the first diagnostic instead.

\textbf{Distillation from strong teachers} (using a large model's correct reasoning chains as SFT data for a smaller model) is expensive but assumed to be necessary for good reasoning.
The corruption results suggest the content of the teacher's chains may matter less than the structural format they provide.

\textbf{Format-only SFT} has not been studied as a deliberate design choice.
The ablation in this paper shows that training exclusively on format---correct step-by-step structure with blank content---yields post-RL accuracy of 76.3\%, close to full-corruption results.

\section{Corruption Study Design}

Five corruption levels are applied to the GSM8K SFT training set:

\begin{center}
\small
\begin{tabular}{lp{5cm}}
\toprule
\textbf{Level} & \textbf{Transformation} \\
\midrule
0 & Original correct chain-of-thought \\
1 & 20\% of intermediate steps replaced with random numbers \\
2 & 50\% of intermediate steps replaced \\
3 & Entire reasoning chain replaced; correct final answer kept \\
4 & Entire reasoning chain and final answer replaced (full noise) \\
\bottomrule
\end{tabular}
\end{center}

For each level, the same base LM is SFT-trained on the corrupted dataset, then RL-trained using GRPO with a binary correctness reward on GSM8K.

\section{Technical Formulation}

The RL phase uses GRPO with the standard RLVR objective.
The reward $r(y, q) = \mathbf{1}[\text{final answer of } y \text{ matches ground truth for } q]$ is binary, operating only on the final numeric answer.
The RL stage has no access to the SFT corruption level and cannot directly evaluate intermediate reasoning steps.

The key diagnostic is the correlation between pre-RL accuracy (SFT model performance on GSM8K) and post-RL accuracy.
A high correlation would indicate that RL merely amplifies what SFT installs; a low correlation indicates RL discovers reasoning independently.

\section{Experimental Findings}

\begin{center}
\small
\begin{tabular}{lrrl}
\toprule
\textbf{Level} & \textbf{SFT acc.} & \textbf{Post-RL} & \textbf{$\Delta$} \\
\midrule
0 (clean) & 68.4\% & 84.7\% & +16.3 \\
1 (20\% corr.) & 61.2\% & 83.9\% & +22.7 \\
2 (50\% corr.) & 52.8\% & 82.4\% & +29.6 \\
3 (chain repl.) & 41.3\% & 80.1\% & +38.8 \\
4 (full noise) & 23.1\% & 72.6\% & +49.5 \\
\bottomrule
\end{tabular}
\end{center}

Post-RL accuracy spans only a 12.1-point range (72.6\%--84.7\%) across corruption levels that span a 45.3-point range in SFT accuracy.
The correlation between SFT accuracy and post-RL accuracy is $r = 0.41$.

Notably, the RL improvement delta \emph{increases} with corruption severity: more corrupted SFT checkpoints benefit more from RL, as if RL has more to fix.

\section{Ablations and Interpretation}

\begin{itemize}[leftmargin=1em,itemsep=2pt]
  \item \textbf{Format-only SFT:} Training on correct step-by-step structure with blank content yields post-RL accuracy of 76.3\%, close to level~3 (80.1\%), confirming that format is the primary SFT contribution.
  \item \textbf{No SFT:} RL without any SFT pre-training yields 61.4\%---substantially lower than any SFT-then-RL variant, confirming that format scaffolding (even from corrupted data) is a necessary precondition.
  \item \textbf{Corruption type:} Shuffled steps (content preserved, order disrupted) yield 81.2\% post-RL; replaced steps (content disrupted) yield 80.1\%, suggesting content is marginally more important than ordering, but both are recoverable.
  \item \textbf{Attention patterns:} At corruption levels~3--4, attention over corrupted intermediate steps decreases during SFT itself, suggesting the model partially ignores corrupted content at the representational level.
\end{itemize}

\section{What the Guarantee Says}

The paper does not argue that SFT trace quality is irrelevant.
The 12.1-point post-RL range confirms that quality still matters---a clean SFT trace yields meaningfully better post-RL performance than full noise.
The claim is that the marginal value of SFT trace quality is much lower than commonly assumed, and that the relationship is non-linear: moderate corruption (level~1--2) causes almost no post-RL degradation, while severe corruption (level~4) causes only a 12-point loss.

The implications are most practical for settings where high-quality SFT data is expensive or unavailable.
The paper suggests that a format-correct but content-imperfect SFT set may be a viable substitute---especially when RL training is planned for a subsequent stage.

\section*{Reference}
\textbf{Reinforcement Learning Unlearns Corrupted SFT Reasoning.}
arXiv:2611.26756. cs.LG, cs.CL. Submitted 19 November 2026.
\url{https://arxiv.org/abs/2611.26756}

\end{document}
